% ===========================================================================
%  Informe técnico — plantilla LaTeX
%
%  Compilar con:   make            (o bien: latexmk -pdf informe.tex)
%  Limpiar:        make clean      (o bien: latexmk -c)
%
%  CÓMO USAR ESTA PLANTILLA
%  ------------------------
%  1. Rellena los datos del trabajo y de los autores en config/datos.tex.
%  2. Escribe dentro de secciones/. Las cajas grises de «Instrucción»
%     dicen qué va en cada sitio, y \ph{...} marca lo que falta.
%  3. Deja las figuras en figuras/, preferiblemente en PDF.
%  4. ANTES DE ENTREGAR, en config/plantilla.tex:
%         \guiafalse          oculta las cajas de instrucción
%         \marcadoresfalse    deja los [Insertar ...] en negro
%     y añade  disable  a las opciones de todonotes para ocultar las
%     notas de pendiente.
%
%  Añadir o quitar secciones se hace aquí, en la lista de \input del cuerpo:
%  este archivo solo ensambla, no contiene texto del informe.
% ===========================================================================

\documentclass[12pt,a4paper,oneside]{article}

\input{config/preambulo}
\input{config/plantilla}
% !TeX root = ../informe.tex
% ===========================================================================
%  Datos de la portada y metadatos del PDF.
%  Este es el único archivo que hay que editar con los datos del trabajo.
% ===========================================================================

\newcommand{\tituloInforme}{Título del informe}
\newcommand{\subtituloInforme}{Subtítulo o alcance del trabajo}
\newcommand{\nombreCurso}{Nombre del curso}
\newcommand{\nombrePrograma}{Nombre del programa}
\newcommand{\nombreUniversidad}{Nombre de la institución}
\newcommand{\tipoTrabajo}{Tipo de trabajo}

% --- Integrantes del equipo ------------------------------------------------
% Un par nombre/correo por integrante. Dejar vacíos los que sobren: la
% portada omite automáticamente los que estén en blanco.
\newcommand{\autorUno}{Nombre del primer autor}
\newcommand{\correoUno}{correo@ejemplo.org}

\newcommand{\autorDos}{}
\newcommand{\correoDos}{}

\newcommand{\autorTres}{}
\newcommand{\correoTres}{}

\newcommand{\autorCuatro}{}
\newcommand{\correoCuatro}{}

\newcommand{\nombreDocente}{Nombre del docente}
\newcommand{\ciudadFecha}{Ciudad, \today}

% --- Metadatos del PDF -----------------------------------------------------
% pdfauthor con varios autores: sepáralos con comas dentro de la clave.
\hypersetup{
    pdftitle  = {\tituloInforme},
    pdfauthor = {\autorUno},
    pdfsubject= {\nombreCurso\ --- \nombrePrograma},
    pdfkeywords = {palabra clave, palabra clave, palabra clave}
}


\begin{document}

% --- Páginas preliminares (numeración romana) ------------------------------
\pagenumbering{roman}

\input{config/portada}

% !TeX root = ../informe.tex
% ===========================================================================
%  Resumen ejecutivo
%  Extensión sugerida: 150-250 palabras, en un solo párrafo.
%  Debe responder: qué se construyó, cómo, y qué resultado se obtuvo.
%  Se escribe al final, cuando el resto del informe ya está listo.
% ===========================================================================

\section*{Resumen}
\addcontentsline{toc}{section}{Resumen}

\ph{Resumen del trabajo: problema, solución y resultado, en un párrafo.}

\vspace{0.5cm}
\noindent\textbf{Palabras clave:} \ph{tres a cinco términos, separados por comas}


\cleardoublepage
\tableofcontents

% Índices de figuras y tablas. Si el informe final no lleva alguno de los
% dos, comenta la línea correspondiente en vez de dejar una página vacía.
\cleardoublepage
\listoffigures

\cleardoublepage
\listoftables

% Índice de pendientes. Útil mientras se escribe; coméntalo al entregar.
\cleardoublepage
\listoftodos[Pendientes]

% --- Cuerpo del informe ----------------------------------------------------
\cleardoublepage
\pagenumbering{arabic}
\setcounter{page}{1}

% !TeX root = ../informe.tex
% ===========================================================================
%  1. Introducción
%  Contenido esperado: contexto del problema, motivación del trabajo,
%  breve descripción de la solución y estructura del documento.
%  Extensión sugerida: 1-2 páginas.
% ===========================================================================

\section{Introducción}
\label{sec:introduccion}

\ph{Contexto y motivación: en qué ámbito se sitúa el trabajo y qué lo hace
pertinente.}

\subsection{Planteamiento del problema}
\label{sec:problema}

\instruccion{Enuncia el problema, no la solución. Si el párrafo ya nombra la
tecnología elegida, es que está describiendo la solución: reescríbelo.}

\ph{Qué necesidad resuelve el sistema y a quién afecta.}

\subsection{Justificación}
\label{sec:justificacion}

\ph{Por qué se aborda con este enfoque y no con otro más simple.}

\subsection{Estructura del documento}
\label{sec:estructura}

\ph{Breve recorrido por las secciones siguientes.}

% !TeX root = ../informe.tex
% ===========================================================================
%  2. Objetivos
%  Deben ser verificables y estar alineados con los criterios de aceptación
%  recogidos en la sección de resultados.
% ===========================================================================

\section{Objetivos}
\label{sec:objetivos}

\subsection{Objetivo general}
\label{sec:objetivo-general}

\instruccion{Un único enunciado, con un verbo comprobable. «Mejorar» o
«optimizar» no se pueden verificar; «implementar», «medir» o «reducir en un
X\,\%», sí.}

\ph{Objetivo general.}

\subsection{Objetivos específicos}
\label{sec:objetivos-especificos}

\instruccion{Cada objetivo específico debería poder rastrearse hasta un
resultado de la \cref{sec:resultados}. Si uno no tiene evidencia asociada,
sobra aquí o falta allí.}

\begin{itemize}
    \item \ph{Objetivo específico 1}
    \item \ph{Objetivo específico 2}
    \item \ph{Objetivo específico 3}
\end{itemize}

% !TeX root = ../informe.tex
% ===========================================================================
%  3. Alcance funcional
%
%  Qué hace el sistema y qué deliberadamente no hace. Es contexto para las
%  secciones de arquitectura, no el capítulo principal: basta con recoger
%  las funcionalidades de forma compacta y señalar lo que se acotó.
% ===========================================================================

\section{Alcance funcional}
\label{sec:alcance}

\subsection{Roles y permisos}
\label{sec:roles}

\instruccion{Una matriz de funcionalidad por rol se lee de un vistazo y
evita párrafos repetitivos. Marca con «Sí» y «No»; si algún permiso depende
de una condición (por ejemplo, actuar solo sobre lo propio), dilo en la
celda. Añade una columna por rol que exista de verdad.}

\begin{table}[H]
  \centering
  \caption{Matriz de roles y permisos.}
  \label{tab:roles}
  \begin{tabular}{L{6.2cm} C{2.6cm} C{2.6cm}}
    \toprule
    \textbf{Funcionalidad} & \ph{Rol 1} & \ph{Rol 2} \\
    \midrule
    \ph{...} & \ph{...} & \ph{...} \\
    \bottomrule
  \end{tabular}
\end{table}

\subsection{Módulos y pantallas}
\label{sec:modulos}

\instruccion{Los módulos funcionales con sus pantallas. Conviene que esta
tabla case con las piezas de la \cref{tab:componentes}: si un módulo no
corresponde a ninguna pieza, o al revés, es señal de que la descomposición
no está clara.}

\begin{table}[H]
  \centering
  \caption{Módulos y pantallas del sistema.}
  \label{tab:modulos}
  \begin{tabular}{L{3.2cm} L{4.4cm} L{5.4cm}}
    \toprule
    \textbf{Módulo} & \textbf{Pantalla} & \textbf{Descripción} \\
    \midrule
    \ph{...} & \ph{...} & \ph{...} \\
    \bottomrule
  \end{tabular}
\end{table}

\subsection{Fuera de alcance}
\label{sec:fuera-alcance}

\instruccion{Enumera lo excluido. No es una disculpa: delimitar el alcance
es una decisión de diseño, y declararla evita que se lea como una carencia.}

\ph{Funcionalidades excluidas y el motivo.}

% !TeX root = ../informe.tex
% ===========================================================================
%  4. Arquitectura de la solución
%
%  Sección central del informe: la estructura del sistema y, sobre todo, su
%  justificación. El modelo de datos va aparte, en la sección 5.
%
%  Deja las figuras en figuras/ con estos nombres, en PDF:
%      contexto.pdf · contenedores.pdf · componentes.pdf · despliegue.pdf
%  Si alguna falta, el documento compila igual y deja un recuadro visible
%  en su lugar.
% ===========================================================================

\section{Arquitectura de la solución}
\label{sec:arquitectura}

\instruccion{Esta sección responde a una sola pregunta: por qué el sistema
está partido así y no de otra manera. Describir los componentes es la parte
fácil; lo que se evalúa es la justificación. Apóyate en las vistas para
mostrar la estructura y en el registro de decisiones para explicarla.}

\subsection{Visión general}
\label{sec:vision-general}

\ph{Dos o tres párrafos: qué estilo arquitectónico se aplica, en cuántas
piezas se divide el sistema y qué principio gobierna el reparto.}

\subsection{Notación}
\label{sec:notacion}

\instruccion{Di en qué notación están los diagramas y con qué herramienta se
generaron. Si salen de una única descripción textual, dilo: es el argumento
de que las vistas son consistentes entre sí por construcción.}

Los diagramas siguen el modelo C4 \cite{ejemplo-web}: contexto, contenedores
y componentes, más una vista de despliegue. \ph{Herramienta y fuente desde
la que se generan.}

\subsection{Vistas arquitectónicas}
\label{sec:vistas}

\instruccion{Una subsección por vista. En cada una: el diagrama, y debajo la
prosa que explica lo que el diagrama no puede decir — por qué esa frontera,
qué pasa si un elemento falla, qué se decidió a propósito. Un diagrama sin
texto no comunica; un texto que solo enumera lo que ya se ve en el diagrama,
tampoco. Quita las vistas que no apliquen en vez de dejarlas vacías.}

\subsubsection{\texorpdfstring{\textsf{V-01}}{V-01} Contexto del sistema}
\label{v:01}

\begin{figure}[H]
  \centering
  \diagrama{contexto}
  \caption{Vista de contexto: el sistema, sus usuarios y los sistemas
           externos con los que se relaciona.}
  \label{fig:contexto}
\end{figure}

\ph{Qué resuelve el sistema y para quién. Un párrafo.}

\subsubsection{\texorpdfstring{\textsf{V-02}}{V-02} Contenedores}
\label{v:02}

\instruccion{Suele ser la vista principal del documento: es donde se ve la
descomposición completa del sistema en piezas ejecutables. Si el diagrama no
cabe legible en el bloque de texto, \texttt{\textbackslash diagramapagina}
lo lleva a página completa.}

\begin{figure}[p]
  \thisfloatpagestyle{empty}
  \centering
  \diagramapagina{contenedores}
  \caption{Vista de contenedores: piezas ejecutables del sistema y sus
           almacenes de datos.}
  \label{fig:contenedores}
\end{figure}

\begin{table}[H]
  \centering
  \caption{Composición del sistema.}
  \label{tab:componentes}
  \begin{tabular}{L{3.4cm} L{6.4cm} L{3.6cm}}
    \toprule
    \textbf{Pieza} & \textbf{Responsabilidad} & \textbf{Tecnología} \\
    \midrule
    \id{\ph{...}} & \ph{...} & \ph{...} \\
    \id{\ph{...}} & \ph{...} & \ph{...} \\
    \id{\ph{...}} & \ph{...} & \ph{...} \\
    \bottomrule
  \end{tabular}
\end{table}

\ph{Recorrido de la tabla: qué agrupa cada pieza y por qué la frontera cae
donde cae. Si alguna no tiene almacén propio, explica de dónde saca los
datos.}

\paragraph*{Comunicación.}
\ph{Qué pieza llama a cuál, con qué protocolo, y qué ocurre si el destino no
responde.}

\subsubsection{\texorpdfstring{\textsf{V-03}}{V-03} Componentes}
\label{v:03}

\instruccion{Detalla una sola pieza, la que concentre la lógica de negocio,
y di explícitamente por qué esa y no otra. Un diagrama de componentes por
cada pieza infla el documento sin añadir nada.}

\begin{figure}[p]
  \thisfloatpagestyle{empty}
  \centering
  \diagramapagina{componentes}
  \caption{Vista de componentes de \id{\ph{pieza}}.}
  \label{fig:componentes}
\end{figure}

\ph{Recorrido de una petición representativa a través de los componentes.}

\subsubsection{\texorpdfstring{\textsf{V-04}}{V-04} Despliegue}
\label{v:04}

\begin{figure}[p]
  \thisfloatpagestyle{empty}
  \centering
  \diagramapagina{despliegue}
  \caption{Vista de despliegue: nodos de ejecución, redes y volúmenes.}
  \label{fig:despliegue}
\end{figure}

\ph{Cómo se materializa la arquitectura en tiempo de ejecución: qué nodos y
redes hay, qué aísla cada una y dónde persisten los datos. El detalle
operativo (comandos, variables de entorno) va en el manual de despliegue, no
aquí: duplicarlo garantiza que los dos se desincronicen.}

\subsection{Seguridad y control de acceso}
\label{sec:seguridad}

\ph{Dónde se autentica el usuario, quién emite la credencial, quién la
valida y cómo llega la identidad a cada pieza del sistema.}

\subsection{Registro de decisiones de arquitectura}
\label{sec:decisiones}

\instruccion{Cada decisión relevante con su alternativa descartada. Una
decisión sin alternativa no es una decisión: es una descripción. Este
registro es lo que distingue un documento de arquitectura de un manual de
usuario del propio sistema. Referencia las decisiones desde el texto con
\texttt{\textbackslash adrref\{01\}}. Repite el bloque una vez por decisión.}

\decision{01}{\ph{Título de la decisión}}
  {\ph{Qué problema o fuerza obligó a decidir.}}
  {\ph{Qué se decidió.}}
  {\ph{Qué otra opción se consideró y por qué se descartó.}}
  {\ph{Qué se gana y qué se paga por ello.}}

\instruccion{Si no sabes por dónde empezar, las decisiones que casi siempre
hay que justificar son: cómo se reparte la persistencia entre piezas; dónde
se validan las reglas de negocio; si el acceso pasa por un punto único o
cada cliente habla directamente con cada servicio; y qué se resuelve en
tiempo de construcción frente a tiempo de ejecución.}

% !TeX root = ../informe.tex
% ===========================================================================
%  5. Modelo de datos
%
%  Lo que se evalúa aquí no es solo que las tablas estén normalizadas, sino
%  que el reparto de los datos entre las piezas del sistema esté razonado.
% ===========================================================================

\section{Modelo de datos}
\label{sec:modelo-datos}

\subsection{Estrategia de persistencia}
\label{sec:estrategia-datos}

\instruccion{Explica el reparto: qué almacén pertenece a qué pieza y qué
principio se sigue. Si se tomó un atajo respecto del diseño ideal — una sola
instancia del motor en lugar de varias, por ejemplo — dilo aquí y
justifícalo: es una pregunta previsible en la defensa.}

\ph{Cómo se reparten los datos, qué implica ese reparto y cómo se garantiza.}

\begin{table}[H]
  \centering
  \caption{Reparto de almacenes de datos.}
  \label{tab:bases}
  \begin{tabular}{L{3.2cm} L{3.4cm} L{6.6cm}}
    \toprule
    \textbf{Almacén} & \textbf{Pieza dueña} & \textbf{Contenido} \\
    \midrule
    \id{\ph{...}} & \id{\ph{...}} & \ph{...} \\
    \id{\ph{...}} & \id{\ph{...}} & \ph{...} \\
    \bottomrule
  \end{tabular}
\end{table}

\subsection{Diagrama entidad-relación}
\label{sec:der}

\instruccion{Un DER por almacén, o uno global con las fronteras marcadas. Lo
segundo comunica mejor la independencia de los datos, que es justo lo que se
evalúa. La figura se espera en \texttt{figuras/modelo-datos.pdf}.}

\begin{figure}[H]
  \centering
  \diagrama{modelo-datos}
  \caption{Modelo entidad-relación, con la frontera de cada almacén.}
  \label{fig:der}
\end{figure}

\subsection{Diccionario de datos}
\label{sec:diccionario}

\instruccion{Una tabla por entidad. No hace falta repetir todos los tipos
SQL: basta el campo, su tipo lógico y qué significa. Los tipos exactos están
en los scripts del anexo. Repite la tabla una vez por entidad relevante.}

\begin{table}[H]
  \centering
  \caption{Entidad \id{\ph{nombre}}.}
  \label{tab:dic-entidad}
  \begin{tabular}{L{3.4cm} L{3.2cm} L{6.6cm}}
    \toprule
    \textbf{Campo} & \textbf{Tipo} & \textbf{Descripción} \\
    \midrule
    \ph{...} & \ph{...} & \ph{...} \\
    \bottomrule
  \end{tabular}
\end{table}

\subsection{Integridad entre almacenes}
\label{sec:integridad}

\instruccion{Es el apartado que más distingue un diseño distribuido de un
monolito con varias tablas. Explica qué referencias cruzan la frontera entre
almacenes, por qué no pueden ser claves foráneas y dónde se valida entonces
la integridad. Si el sistema tiene un solo almacén, sustituye este apartado
por las restricciones de integridad referencial y dilo.}

\ph{Referencias lógicas sin clave foránea: cuáles son, dónde se validan y
qué ocurre si la validación falla.}

\subsection{Restricciones que implementan reglas de negocio}
\label{sec:restricciones}

\instruccion{Qué reglas de la \cref{sec:reglas-negocio} están sostenidas por
el propio motor de base de datos, y por qué ahí y no en el código. Es un
argumento de concurrencia: una validación en la aplicación se puede saltar si
dos peticiones llegan a la vez. Aquí va solo el mecanismo; el enunciado de la
regla vive en la sección 6.}

\ph{Restricciones e índices que sostienen las reglas y las consultas más
frecuentes.}

% !TeX root = ../informe.tex
% ===========================================================================
%  6. Reglas de negocio
%
%  Lo que se evalúa es la trazabilidad: cada regla, dónde se implementa y
%  cómo se comprueba.
%
%  Los \label{rn:NN} de la tabla son los que hacen funcionar \rn{NN} desde
%  el resto del documento: si añades filas, añade también su \label.
% ===========================================================================

\section{Reglas de negocio}
\label{sec:reglas-negocio}

\subsection{Catálogo de reglas}
\label{sec:catalogo-reglas}

\instruccion{Una fila por regla, con la columna de implementación rellenada
de verdad: «\id{ms-pedidos}, \id{PedidoService}» dice algo; «se valida en el
backend» no. Esta tabla es la que se mira para puntuar este criterio.}

\begin{table}[H]
  \centering
  \caption{Reglas de negocio y dónde se implementa cada una.}
  \label{tab:reglas}
  \begin{tabular}{C{1.6cm} L{7.2cm} L{4.4cm}}
    \toprule
    \textbf{ID} & \textbf{Descripción} & \textbf{Implementación} \\
    \midrule
    \textbf{RN-01}\label{rn:01} & \ph{Enunciado de la regla} & \ph{...} \\
    \textbf{RN-02}\label{rn:02} & \ph{Enunciado de la regla} & \ph{...} \\
    \textbf{RN-03}\label{rn:03} & \ph{Enunciado de la regla} & \ph{...} \\
    \bottomrule
  \end{tabular}
\end{table}

\subsection{Regla crítica}
\label{sec:regla-critica}

\instruccion{Elige la regla que más riesgo concentra — típicamente una que
falla solo bajo concurrencia o con datos límite — y dedícale espacio: qué
pasa si no se cumple, dónde se resuelve y por qué ahí. Renombra el apartado
con el nombre de la regla.}

\ph{Mecanismo elegido para \rn{01} y por qué resiste el caso adverso.}

\subsection{Máquina de estados}
\label{sec:estados}

\instruccion{Si alguna entidad del sistema pasa por estados, este es su
sitio. Un diagrama de estados pequeño comunica esto mejor que un párrafo; si
no hay estados que modelar, borra el apartado en vez de dejarlo vacío.}

\ph{Estados y transiciones permitidas.}

\subsection{Traducción de reglas a respuestas de la API}
\label{sec:errores}

\instruccion{Qué código devuelve la API cuando se viola cada regla. Importa
más de lo que parece: un 500 genérico ante una violación previsible es un
defecto, y un código específico con un mensaje claro es lo que permite que la
interfaz reaccione bien.}

\begin{table}[H]
  \centering
  \caption{Correspondencia entre violación de regla y respuesta HTTP.}
  \label{tab:errores}
  \begin{tabular}{C{1.6cm} C{2.4cm} L{8.6cm}}
    \toprule
    \textbf{Regla} & \textbf{HTTP} & \textbf{Código de error y significado} \\
    \midrule
    \ph{...} & \ph{...} & \ph{...} \\
    \bottomrule
  \end{tabular}
\end{table}

% !TeX root = ../informe.tex
% ===========================================================================
%  7. Resultados
%
%  Cierra el círculo con los objetivos (sección 2) y con los criterios de
%  aceptación acordados. Los \label{ca:N} son los que hacen funcionar
%  \crit{N} desde el resto del documento: si añades filas, añade su \label.
% ===========================================================================

\section{Resultados}
\label{sec:resultados}

\subsection{Funcionalidades implementadas}
\label{sec:funcionalidades}

\instruccion{Qué quedó operativo, con capturas de las pantallas principales.
Dos o tres capturas bien elegidas valen más que ocho: las que muestran el
flujo central del sistema. El resto va al anexo de capturas.}

\ph{Resumen de lo implementado.}

\subsection{Cumplimiento de los criterios de aceptación}
\label{sec:cumplimiento}

\instruccion{Un criterio por fila, con la evidencia que respalda cada «Sí».
Es la tabla que un evaluador busca primero. Si alguno no se cumple,
decláralo: un criterio marcado como cumplido que luego falla en la demo
cuesta más que uno declarado pendiente.}

\begin{table}[H]
  \centering
  \caption{Cumplimiento de los criterios de aceptación.}
  \label{tab:cumplimiento}
  \begin{tabularx}{\textwidth}{C{1.4cm} Y C{2cm} L{3cm}}
    \toprule
    \textbf{\#} & \textbf{Criterio} & \textbf{Estado} & \textbf{Evidencia} \\
    \midrule
    \textsf{CA-1}\label{ca:1} & \ph{Criterio acordado} & \ph{...} & \ph{...} \\
    \textsf{CA-2}\label{ca:2} & \ph{Criterio acordado} & \ph{...} & \ph{...} \\
    \textsf{CA-3}\label{ca:3} & \ph{Criterio acordado} & \ph{...} & \ph{...} \\
    \bottomrule
  \end{tabularx}
\end{table}

\subsection{Verificación de las reglas de negocio}
\label{sec:verificacion-reglas}

\instruccion{Cómo se comprobó que cada regla de la \cref{tab:reglas}
funciona. No hace falta un capítulo de pruebas: una tabla con el escenario,
el resultado esperado y el obtenido es suficiente y se lee en treinta
segundos. La colección de peticiones va al anexo.}

\begin{table}[H]
  \centering
  \caption{Comprobación de las reglas de negocio.}
  \label{tab:pruebas-reglas}
  \begin{tabular}{C{1.6cm} L{6.6cm} L{4.8cm}}
    \toprule
    \textbf{Regla} & \textbf{Escenario} & \textbf{Resultado} \\
    \midrule
    \ph{...} & \ph{...} & \ph{...} \\
    \bottomrule
  \end{tabular}
\end{table}

\subsection{Limitaciones}
\label{sec:limitaciones}

\instruccion{Qué no se alcanzó y por qué. Declararlo suma: demuestra que se
conoce el propio sistema. Omitir una limitación que el evaluador descubre en
la demo resta el doble.}

\ph{Limitaciones conocidas.}

% !TeX root = ../informe.tex
% ===========================================================================
%  8. Conclusiones
%  Conclusiones sobre los objetivos planteados, no un resumen del trabajo
%  realizado. Cada conclusión debería poder rastrearse hasta un resultado.
% ===========================================================================

\section{Conclusiones}
\label{sec:conclusiones}

\instruccion{Una conclusión por objetivo específico de la
\cref{sec:objetivos-especificos}, apoyada en la evidencia de la
\cref{sec:resultados}. Si una conclusión no se puede rastrear hasta un
resultado, es una opinión.}

\ph{Conclusiones respecto de los objetivos.}

\subsection{Lecciones aprendidas}
\label{sec:lecciones}

\ph{Qué se haría distinto con lo aprendido.}

\subsection{Trabajo futuro}
\label{sec:trabajo-futuro}

\ph{Extensiones naturales del sistema, incluidas las excluidas del alcance
en la \cref{sec:fuera-alcance}.}


% --- Referencias -----------------------------------------------------------
\cleardoublepage
\printbibliography[heading=bibintoc,title={Referencias}]

% --- Anexos ----------------------------------------------------------------
\appendix
% !TeX root = ../informe.tex
% ===========================================================================
%  Anexos
%  Material de respaldo: scripts completos, colecciones de pruebas, capturas.
%  Con \appendix ya activo, las secciones se numeran A, B, C...
%  Borra los anexos que no apliquen: un anexo vacío resta.
% ===========================================================================

\section{Scripts de base de datos}
\label{anexo:sql}

\ph{DDL completo y datos de prueba, o la ruta del repositorio donde viven.}

\section{Colección de pruebas de la API}
\label{anexo:api}

\ph{Peticiones de ejemplo, agrupadas por servicio.}

\section{Capturas de la aplicación}
\label{anexo:capturas}

\ph{Pantallas de los flujos principales.}

\section{Repositorio del proyecto}
\label{anexo:repositorio}

\ph{URL del repositorio e instrucciones de acceso.}


\end{document}
